\section{Referencias}

\begin{thebibliography}{99}

\bibitem{covenin3172}
Comisión Venezolana de Normas Industriales (COVENIN), \textit{Norma Venezolana COVENIN 3172-95: Transformadores de Potencia. Métodos de Ensayo}, FONDONORMA, Caracas, Venezuela, 1995.

\bibitem{covenin536}
Comisión Venezolana de Normas Industriales (COVENIN), \textit{Norma Venezolana COVENIN 536-94: Transformadores de Potencia. Generalidades}, FONDONORMA, Caracas, Venezuela, 1994.

\bibitem{ieee5712}
IEEE Power Engineering Society, \textit{ANSI/IEEE C57.12.91-2001: IEEE Standard Test Code for Dry-Type Distribution and Power Transformers}, Institute of Electrical and Electronics Engineers, Nueva York, 2001.

\bibitem{chapman}
S. J. Chapman, \textit{Máquinas Eléctricas}, 5.ª ed. México D.F.: McGraw-Hill / Interamericana Editores, 2012.

\bibitem{fraile}
J. Fraile Mora, \textit{Máquinas Eléctricas}, 6.ª ed. Madrid, España: McGraw-Hill, 2008.

\bibitem{mit}
E. E. Staff del M.I.T., \textit{Circuitos Magnéticos y Transformadores}, Barcelona, España: Editorial Reverté, 1981.

\bibitem{ras}
E. Ras, \textit{Transformadores de Potencia, de Medida y de Protección}, 4.ª ed. Barcelona, España: Marcombo Boixareu Editores, 1978.

\bibitem{guia4}
N. Ojeda y J. Pérez, \textit{Práctica Nº 4: Transformador Monofásico. Modelo}, Laboratorio de Máquinas Eléctricas I, Departamento de Potencia, Escuela de Ingeniería Eléctrica, Facultad de Ingeniería, Universidad Central de Venezuela, Caracas, 2012.

\bibitem{guia6}
Universidad Central de Venezuela, \textit{Práctica Tema III: Laboratorio de Máquinas I. Transformador Monofásico: Rendimiento}, Escuela de Ingeniería Eléctrica, Departamento de Potencia, Caracas, Venezuela.

\end{thebibliography}
